% =================================================
% Домашняя работа: ученическая часть — скорректированная версия
% =================================================
\homeworktitle
  {Домашняя работа}
  {Степени логарифма, модуль и сведение к квадратному неравенству}
\studentline

\begin{routebox}[Маршрут]
  Для полного решения придерживайтесь порядка:
  \begin{enumerate}
    \item прочитайте нотацию и найдите ОДЗ;
    \item преобразуйте чётные степени с помощью модуля;
    \item введите \(y=\log_b|f(x)|\), затем \(t=y^2\ge0\);
    \item решите квадратное неравенство;
    \item вернитесь к логарифму и к исходной переменной.
  \end{enumerate}
\end{routebox}

\begin{homeworktask}[Три похожие записи]{1}
  а) Расшифруйте:
  \[
    \lg^3 A=\answerblank[35mm],
    \qquad
    \lg(A^3)=\answerblank[35mm],
  \]
  \[
    \log_3 A=\answerblank[70mm].
  \]

  б) Вычислите и сравните:
  \[
    \lg^2 1000=\answerblank[25mm],
    \qquad
    \lg(1000^2)=\answerblank[25mm].
  \]

  в) Объясните, почему \(\lg^2 A\) и \(\lg(A^2)\) не являются одной и той же записью.
  \answerspace{22mm}
\end{homeworktask}

\newpage
\begin{homeworktask}[Исправьте формулу]{2}
  Ученик записал:
  \[
    \lg((x-1)^2)=2\lg(x-1).
  \]
  \begin{enumerate}
    \item На какой области определена левая часть?
    \item На какой области определена правая часть?
    \item Запишите равенство, верное на всей области определения левой части.
  \end{enumerate}
  \answerspace{36mm}
\end{homeworktask}

\begin{homeworktask}[Алгебраическое ядро]{3}
  Решите неравенство
  \[
    y^4-y^2-12\ge0.
  \]
  Введите \(t=y^2\) и обязательно укажите ограничение на \(t\).
  \[
    t^2-t-12=(t-\answerblank[12mm])(t+\answerblank[12mm]).
  \]
  \answerspace{35mm}
  \shortanswerline
\end{homeworktask}

\newpage
\begin{homeworktask}[Двойное неравенство с модулем]{4}
  Решите два неравенства. В первом случае сначала положите \(u=x-2\),
  во втором --- \(u=x^2-4\).
  \[
    \text{а) }0<|x-2|\le3,
    \qquad
    \text{б) }0<|x^2-4|\le1.
  \]
  Для каждого задания изобразите ответ на числовой прямой.
  \answerspace{62mm}
\end{homeworktask}

\newpage
\begin{homeworktask}[Полное решение с опорой]{5}
  Решите неравенство
  \[
    \log_2^4((x-1)^2)-\log_2^2((x-1)^4)\ge192.
  \]

  \structuredsolutionmap

  \begin{hintbox}[Контрольные ориентиры]
    Сначала найдите ОДЗ. Затем выберите выражение, которое удобно обозначить
    через \(y\), чтобы исходное неравенство превратилось в биквадратное.
  \end{hintbox}

  \answerspace{42mm}
\end{homeworktask}

\begin{homeworktask}[Выберите корректную замену]{6}
  Для выражения
  \[
    \lg^4((x^2-1)^2)-\lg^2((x^2-1)^4)
  \]
  предложены замены:
  \[
    \text{A. }t=\lg^2|x^2-1|,
    \qquad
    \text{Б. }t=\log_2|x^2-1|,
  \]
  \[
    \text{В. }t=\lg(|x^2-1|^2).
  \]
  Какая замена непосредственно превращает выражение в квадратный многочлен
  относительно \(t\)? Расшифруйте выбранную замену без сокращённой записи.
  \answerspace{36mm}
\end{homeworktask}

\newpage
\begin{homeworktask}[Самостоятельное решение]{7}
  Решите неравенство
  \[
    \log_3^4((x+2)^2)-\log_3^2((x+2)^4)\ge192.
  \]
  Запишите полное решение с ОДЗ и точным ответом.
  \reducedsolutionmap
  \answerspace{78mm}
\end{homeworktask}

\newpage
\begin{homeworktask}[Полностью самостоятельное решение]{8}
  Решите неравенство
  \[
    \log_2^4((x^2-1)^2)-\log_2^2((x^2-1)^4)\ge192.
  \]
  Запишите полное решение с ОДЗ и точным ответом. Карта решения и
  контрольные ориентиры к этому заданию не даны.
  \answerspace{85mm}
\end{homeworktask}

\newpage
\begin{homeworktask}[Обобщение]{9}
  Пусть \(b>1\), \(A\ne0\). Решите относительно \(A\) неравенство
  \[
    \log_b^4(A^2)-\log_b^2(A^4)\ge192.
  \]
  Ответ выразите через степени основания \(b\).
  \answerspace{48mm}
  \shortanswerline
\end{homeworktask}

\begin{selfcheck}
  \begin{itemize}
    \item Я различаю степень значения логарифма, степень аргумента и основание.
    \item При преобразовании \(\log_b(z^{2n})\) сохранил \(|z|\).
    \item Для второй замены указал \(t\ge0\).
    \item После неравенства с квадратом сохранил две ветви.
    \item В ветви малого аргумента сохранил строгое условие аргумент \(>0\).
  \end{itemize}
\end{selfcheck}
